\documentclass[12pt]{article}
\usepackage[utf8]{inputenc}
\usepackage[T1]{fontenc}
\usepackage[margin=1in]{geometry}
\usepackage{setspace}
\usepackage{xcolor}
\usepackage[colorlinks=true, urlcolor=blue!60!black, linkcolor=blue!60!black]{hyperref}
\setlength{\parindent}{0em}
\setlength{\parskip}{0.6em}
\renewcommand{\baselinestretch}{1.0}
\begin{document}
\begin{center}
{\LARGE \textit{Fletcher v. Peck}, 10 U.S. 87 (1810)\par}
\vspace{0.5em}
{\large Case Brief --- Two-Page Edition\par}
\end{center}

\section{Memory Jogger}
In the first decision striking down a state statute, the Court held that Georgia's Yazoo land grant was a ``contract'' under the Contract Clause that Georgia could not rescind against innocent purchasers, despite the bribery behind it (\href{https://www.courtlistener.com/opinion/84935/fletcher-v-peck/}{10 U.S. 87}).

\section{Detailed Case Facts}
In 1795 a bribed Georgia legislature sold about 35 million acres along the Yazoo River to land companies for roughly 1.5 cents per acre. A newly elected legislature declared the grant void in 1796, but tracts had already been resold to innocent purchasers. Peck conveyed 15,000 acres to Fletcher with covenants that title was good, and Fletcher sued for breach in what Justice Johnson saw as a ``feigned'' test case.

\section{Procedural History}
The federal circuit court in Massachusetts gave judgment for Peck on all four counts; Fletcher brought a writ of error, and the Supreme Court affirmed in 1810.

\section{Judicial Votes}
Affirmed, 4--1 on the reasoning. Marshall wrote for Washington, Livingston, and Todd; Johnson concurred in part and dissented in part; Cushing and Chase did not participate.

\section{Holding}
A completed state land grant is an executed contract under Article~I, \S~10, so a later legislature cannot rescind it once title has vested in bona fide purchasers, and courts in private suits will not void statutes by probing legislators' corrupt motives.

\section{Analysis of Opinions}
Marshall reasoned that courts cannot feasibly inquire into individual legislators' motives and that equity protects purchasers without notice. ``[W]hen a law is in its nature a contract, when absolute rights have vested under that contract, a repeal of the law cannot devest those rights.'' Because a grant extinguishes the grantor's right, and the Clause's text draws no line between private and public contracts, Georgia was bound. Indian title was a right of occupancy compatible with Georgia's fee simple. Johnson agreed Georgia could not revoke its grant but rested on natural law, ``a principle which will impose laws even on the deity,'' warning that the Contract Clause reading would threaten eminent domain and other ordinary legislation; he dissented on Indian title.

\section{Examples: Future Applications}
\textbf{Same-Side:} a legislature converting a privately chartered university into a state institution impairs an executed contract, as in \textit{Dartmouth College}. \textbf{Opposite-Side:} prospective repeal of an annual solar subsidy that created no vested rights is valid. \textbf{Fence-Sitter:} a state suing in equity to cancel a bribe-procured mineral lease before any transfer to innocent parties, a question Marshall left open.

\section{Critique}
Originalists and textualists argue the Clause targeted private debtor-relief laws and that an executed grant is no ongoing ``obligation of contracts,'' and they fault the Court for hearing a collusive suit. Progressive historians see a bribed giveaway shielded from democratic reversal, and Indian law scholars see a step toward \textit{Johnson v. M'Intosh}.

\section{Key Quotations}
\begin{quote}
``A contract executed, as well as one which is executory, contains obligations binding on the parties. A grant, in its own nature, amounts to an extinguishment of the right of the grantor, and implies a contract not to reassert that right. A party is, therefore, always estopped by his own grant.''
\end{quote}
\noindent --- Chief Justice Marshall, Opinion of the Court, 10 U.S. at 137.

\end{document}
